\documentclass[11pt,a4paper]{article}

% ---------------------------------------------------------------------------
%  Autonomous Robots -- Modularbeit (OTH Amberg-Weiden)
%  Technical documentation of the maze-navigation approach.
%  Compile with: pdflatex Technical_Report.tex   (run twice for references)
% ---------------------------------------------------------------------------

\usepackage[utf8]{inputenc}
\usepackage[T1]{fontenc}
\usepackage[margin=2.3cm]{geometry}
\usepackage{amsmath,amssymb}
\usepackage{booktabs}
\usepackage{array}
\usepackage{longtable}
\usepackage{enumitem}
\usepackage{graphicx}
\usepackage{xcolor}
\usepackage{titlesec}
\usepackage[hidelinks]{hyperref}
\usepackage{caption}

\definecolor{codegray}{gray}{0.95}
\newcommand{\code}[1]{\colorbox{codegray}{\texttt{\small #1}}}

\titleformat{\section}{\normalfont\large\bfseries}{\thesection}{0.6em}{}
\titleformat{\subsection}{\normalfont\normalsize\bfseries}{\thesubsection}{0.6em}{}
\setlist{topsep=2pt,itemsep=1pt,parsep=1pt}

\title{\vspace{-1.2cm}\textbf{Autonomous Maze Navigation of a ROSbot in Webots}\\[0.3em]
\large A Common Controller Architecture with Per-Map Configuration}
\author{Autonomous Robots -- Modularbeit, OTH Amberg-Weiden}
\date{\today}

\begin{document}
\maketitle
\vspace{-0.8cm}

% ---------------------------------------------------------------------------
\begin{abstract}
\noindent This document specifies the navigation approach adopted for the
Autonomous Robots Modularbeit. The objective of every maze world is identical:
a ROSbot must reach the \textbf{blue} pillar first and the \textbf{yellow}
pillar second, in minimum simulation time, without ever touching a wall or
crossing the \textbf{green} poison floor, using no Supervisor and no edits to
the robot or world. The approach reuses a single, algorithmically fixed
navigation controller across all five maze worlds; only the geometry and
tuning constants held in each package's \code{config.py} differ between mazes.
The pipeline couples lidar log-odds SLAM, a persistent depth-obstacle layer for
lidar-blind hazards, HSV-based visual target acquisition, frontier-driven
exploration, A* global planning on a centre-seeking costmap, and a Dynamic
Window Approach (DWA) local planner that keeps wall clearance immune to pose
drift. The algorithmic pattern is adapted from a public reference
implementation, which is credited explicitly to preserve academic integrity.
\end{abstract}

% ---------------------------------------------------------------------------
\section{Problem Statement and Design Principle}

The task defines a fixed mission over an unknown static maze: navigate from the
start pose to the blue pillar, then to the yellow pillar, in the least
simulation time, subject to two hard safety constraints --- no wall contact and
no traversal of the green poison patch. The pillar positions are not provided;
the target must be discovered by on-board vision during the run.

The guiding design principle of the approach is \emph{one algorithm, many
configurations}. Rather than authoring a separate controller per maze, the
project maintains a single navigation stack whose control logic, state machine,
sensor interface and observability are identical everywhere. Each maze world
binds to its own package copy so that the controller remains self-contained and
runs standalone in Webots, but the only substantive difference between those
copies is the set of tuning constants in \code{config.py}. This separation
keeps the algorithm auditable and regression-tested once, while allowing the
geometry of each arena --- grid resolution, arena extent, hazard height bands,
sensor complement --- to be expressed declaratively as data. The class name
(\code{NavigationController}), the Webots device layer
(\code{RobotInterface}) and the observability layer are shared verbatim across
every maze.

% ---------------------------------------------------------------------------
\section{System Architecture}

The controller is decomposed into single-responsibility modules that are
executed once per simulation tick. The module layout is uniform across mazes:

\begin{center}
\small
\begin{tabular}{@{}p{3.2cm}p{12.0cm}@{}}
\toprule
\textbf{Module} & \textbf{Responsibility} \\
\midrule
\code{robot\_io.py} & \code{RobotInterface} --- the only Webots-facing module: device discovery, guarded reads of motors, lidar, RGB and depth cameras, and IR range-finders. \\
\code{odometry.py} & Wheel-encoder and IMU-yaw fusion for the incremental pose estimate. \\
\code{mapping.py} & Log-odds occupancy grid, scan-matching, costmap generation, and the poison and auxiliary floating-wall layers. \\
\code{depth\_model.py} & Depth image to thin-scan conversion (per-column hit/clear) for the floating-wall layer. \\
\code{perception.py} & HSV detection of the blue/yellow pillars and green-poison projection with a forward reflex. \\
\code{frontier.py} / \code{astar.py} & Frontier extraction and clustering; A* global planning with clearance-aware path simplification. \\
\code{local\_planner.py} & Arc-length pure-pursuit carrot plus a Dynamic Window Approach local planner. \\
\code{mak\_02\_controller.py} & \code{NavigationController} --- the mission finite-state machine and the main sensing/planning/driving loop. \\
\code{observability.py} & Leveled logging, a structured JSONL event log, and per-stage fault barriers (vendored from \code{common/}). \\
\code{config.py} & Single source of truth for every tuning constant; the only file that differs per maze. \\
\bottomrule
\end{tabular}
\end{center}

A separate \code{common/} package sits outside every maze folder and is imported
by none of the running controllers, so it cannot influence the code that runs
in Webots. It provides a uniform \code{MazeControllerWrapper} that resolves any
maze identifier to its \code{NavigationController} and exposes only the shared
method contract, together with the canonical \code{observability.py} that each
controller vendors a copy of.

\subsection{Per-tick data flow}

Each control cycle proceeds through the following stages, each wrapped in a
fault barrier so that a single bad frame is logged and skipped rather than
crashing the mission:

\begin{enumerate}
\item \textbf{Sensing.} Read encoders, IMU, lidar scan, RGB and depth images,
and IR ranges through \code{RobotInterface}.
\item \textbf{State estimation.} Advance odometry, then refine the pose by
scan-matching the current lidar scan against the accumulated map.
\item \textbf{Mapping.} Integrate the lidar scan into the log-odds grid,
accumulate depth hits into the auxiliary layer, and project detected poison.
\item \textbf{Perception.} Detect the target pillar and estimate its range.
\item \textbf{Costmap.} Rebuild the inflated, centre-seeking costmap from the
occupancy, auxiliary and poison layers.
\item \textbf{Goal selection.} Choose a frontier goal while exploring, or the
pillar standoff once the target is visually fixed and reachable.
\item \textbf{Planning.} Plan a global A* path to the current goal.
\item \textbf{Driving.} Convert the path to a pure-pursuit carrot and let DWA
select the executed velocity command from the live lidar.
\end{enumerate}

% ---------------------------------------------------------------------------
\section{State Estimation and Mapping}

\subsection{Pose estimation}
The pose is maintained by fusing wheel-encoder odometry with IMU yaw. Because
integrated odometry drifts, the estimate is corrected each tick by
scan-matching the incoming lidar scan against the current occupancy map; the
matched correction keeps the map self-consistent and, critically, keeps wall
clearance referenced to what the robot currently senses rather than to an
accumulating error.

\subsection{Log-odds occupancy grid}
The environment is represented as a 2-D occupancy grid updated in log-odds form
from the single-plane lidar. Cells struck by a beam gain occupancy evidence and
cells traversed by a beam lose it, so transient noise is averaged out over
successive scans. The grid is square and centred on the start pose; its
resolution and extent are maze-specific constants (Section~\ref{sec:config}).

\subsection{Poison layer}
The green poison floor is a lethal constraint. Green pixels detected by the RGB
camera are projected into the grid as a dedicated poison layer that the planner
treats as impassable. A complementary forward reflex
(\code{GREEN\_REFLEX}) blocks forward motion whenever a cluster of projected
green points appears in a narrow corridor directly ahead in the body frame,
guarding against a plan that has not yet caught up with a freshly seen patch.

\subsection{Auxiliary floating-wall layer}
The single-plane 2-D lidar sits at approximately $z=0.10$~m and is therefore
blind to walls that float above, tilt across, or sit below that plane; such
walls are marked \emph{free} by the lidar and would otherwise be driven into.
The depth camera can see them, but only beyond its minimum range of roughly
$0.6$~m, which leaves a ``dead shell'' closer than $0.6$~m where a floating wall
is invisible to every sensor. The only defence in that shell is memory.

The approach therefore back-projects every depth pixel that falls inside the
collision height band $[\,\code{AUX\_Z\_MIN},\,\code{AUX\_Z\_MAX}=\text{robot
height}\,]$ and accumulates per-cell confidence wherever the lidar is blind. A
cell becomes lethal ``aux'' once it reaches \code{AUX\_MIN\_HITS}; a clear
sight-line only \emph{decays} that confidence, and the decay ray starts at the
depth minimum range so a wall already inside the dead shell is never forgotten
during the approach. A lidar-blind gate prevents ordinary full-height walls
from being duplicated into the aux layer, and the confidence saturates at
\code{AUX\_HIT\_CAP} so a confirmed wall survives the few stray clear frames of
a close approach. Panels taller than the robot (pass-under bridges) fall outside
the height band and stay drivable. Aux cells feed both the A* costmap (as walls)
and the DWA obstacle set.

% ---------------------------------------------------------------------------
\section{Perception and Target Acquisition}

Pillars are acquired visually rather than from any hardcoded coordinate. The RGB
image is thresholded in HSV space for the blue and yellow pillar hues; the
detected blob, combined with the known pillar height, yields a bearing and a
range estimate that place the target in the map frame. A running mean over the
most recent detections stabilises the estimate against single-frame noise, and
successive observations from different viewpoints sharpen it. No world
coordinate of any pillar is present in the runtime path; the tabulated
positions in the controller documentation are provided only for verification.

% ---------------------------------------------------------------------------
\section{Exploration, Planning and Control}

\subsection{Visual-frontier exploration}
While the target pillar has not yet been localised, the robot explores the
unknown space by frontier search. Frontiers --- the boundaries between
known-free and unknown cells --- are extracted and clustered, and each candidate
is scored by a utility that rewards information gain and penalises path length,
biased to keep the robot moving forward. The chosen frontier becomes the A*
goal.

Exploration and goal acquisition are blended: the instant the RGB camera fixes
the target pillar \emph{and} A* finds a path to its standoff, the state machine
abandons exploration and commits directly to the pillar
(\code{\_maybe\_go}). This ``visual frontier plus DWA'' blend avoids exhaustive
mapping --- the maze is only explored as far as is necessary to see the next
target.

\subsection{A* on a centre-seeking costmap}
The global planner is an 8-connected, weighted A* over a costmap derived from
the occupancy, auxiliary and poison layers. Two properties matter for the tight
maze gaps. First, obstacle inflation is footprint-accurate: the lethal radius
\code{HARD\_OBS\_DIST} equals the true passage half-width of the ROSbot
($0.116$~m), so A* never routes the wheels into a wall while still leaving a
multi-cell free band down the centre of a corridor. Second, cost decreases
smoothly with clearance out to \code{CENTER\_PREF\_RANGE}, so the planned path
rides the medial axis of every corridor. Path simplification is
clearance-aware: it shortcuts only through genuinely open cells and never
straightens a centred path back against a tight wall. Poison cells are weighted
rather than merely forbidden, so the planner skirts them but retains a route if
one is unavoidable.

\subsection{Pure-pursuit carrot and DWA local planner}
The A* path only decides \emph{where} to aim. The executed velocity command is
chosen every tick by a Dynamic Window Approach that samples feasible
$(v,\omega)$ pairs, rolls each out over a short horizon against the
\emph{live} lidar obstacles, and rejects any rollout passing closer than
\code{DWA\_SAFE\_RADIUS}. Surviving candidates are scored by a weighted sum of
heading alignment to a pure-pursuit carrot, proximity to that carrot, and
clearance (which doubles as corridor centring). Because clearance is evaluated
from the current scan rather than the map, wall avoidance is immune to
accumulated pose drift. A single floored tight-gap speed factor eases the robot
to roughly $0.16$--$0.26$~m/s only where a side wall is genuinely close,
preserving cruising speed elsewhere. The four-wheel skid-steer ROSbot is
commanded as a two-wheel differential drive, mapping $(v,\omega)$ to a left and
a right wheel pair; centreline precision comes from the planner, not from any
drivetrain modification.

\subsection{Mission state machine and recovery}
The mission is governed by a finite-state machine:
\begin{align*}
\text{INIT\_SCAN} \rightarrow \text{EXPLORE\_BLUE} &\rightarrow \text{GO\_BLUE}
\rightarrow \text{EXPLORE\_YELLOW} \\
&\rightarrow \text{GO\_YELLOW} \rightarrow \text{DONE.}
\end{align*}
A \textbf{RECOVERY} state is reachable from any driving state when forward
progress stalls. A windowed-displacement watchdog detects a genuine lack of
progress (rather than mere pose jitter) and triggers recovery, which reverses
when there is room behind and otherwise spins in place, then re-plans. Repeated
recoveries escalate, and an area that cannot be traversed is blacklisted so the
planner tries an alternative frontier.

% ---------------------------------------------------------------------------
\section{Per-Map Configuration}
\label{sec:config}

The five maze worlds run the same algorithm; the differences are confined to
\code{config.py}. The table below summarises the substantive per-map settings
and the reasoning behind each.

\begin{center}
\small
\begin{tabular}{@{}p{1.1cm}p{1.3cm}p{2.5cm}p{7.6cm}@{}}
\toprule
\textbf{Maze} & \textbf{Grid res.} & \textbf{Sensor set} & \textbf{Rationale for the configuration} \\
\midrule
Maze1 & $0.025$~m & lidar + depth + IR & Full stack; floating-wall avoidance designed here, aux band from $z=0.01$~m for near-floor panels. \\
Maze2 & $0.025$~m & lidar + depth + IR & High-resolution SLAM on a $9.0$~m grid; aux band raised to $z=0.03$~m for the floating/tilted blue-approach cluster. \\
Maze3 & $0.025$~m & lidar + depth + IR & Map-consistent SLAM fused with map-based poison on an $11.0$~m grid. \\
Maze4 & $0.04$~m & lidar + depth + IR & Reference tuning; coarser grid on an $11.0$~m arena, tighter aux max-range. \\
Maze5 & $0.05$~m & lidar only & Baseline; no floating walls, so the depth-aux and IR layers are unnecessary. \\
\bottomrule
\end{tabular}
\end{center}

Grid resolution trades runtime against wall sharpness: the tighter arenas with
lidar-blind hazards (Maze1--Maze3) use $0.025$~m for crisp planning, whereas
the open Maze5 baseline uses $0.05$~m. The auxiliary-layer height band is set
per maze from the actual $z$-spans of that world's floating, tilted and
near-floor panels, and the pillar height constant (0.40~m in Maze2, 0.30~m
elsewhere) feeds the vision range estimate. Every geometry constant is read
from the world file and documented in the respective controller's README, but
no world coordinate is hardcoded in the runtime path.

% ---------------------------------------------------------------------------
\section{Observability and Validation}

The controller carries an enterprise-style observability layer. Operator
messages are emitted through leveled, timestamped logging, and a machine-
readable trail is written to \code{maps/run\_events.jsonl}: a run-start record
with the full device inventory, mission state transitions, pillar-reached
splits with timing, recovery entries, status heartbeats, fault records with
tracebacks, the final timing table, and a run-end record. Every control-loop
stage runs inside a \code{guarded\_stage} barrier, individual device reads are
guarded, and a top-level safety net guarantees the wheels are stopped if the
controller ever crashes.

The stack is validated without a simulator. Each package ships a
\code{selftest.py} (unit checks over mapping, scan-matching, A*, frontier
extraction, DWA, perception, the depth-aux mark/clear logic, and IR lookup,
plus a closed-loop ``no-crawl'' regression) and, where applicable, a
\code{dryrun.py} that drives the real carrot-plus-DWA controller through
synthetic straight, narrow-gap, offset, L-corner and pivot corridors and
asserts that the robot moves at cruising speed, stays clear of walls, and
reaches the goal. A \code{smoke\_test.py} in \code{common/} confirms that all
five mazes resolve and satisfy the shared contract. These Webots-free gates are
run after every change, so a tuning adjustment for one maze cannot silently
regress the shared algorithm.

% ---------------------------------------------------------------------------
\section{Attribution and Academic Integrity}

The navigation pattern --- lidar SLAM, frontier-based exploration, and a
\code{move\_base}/\code{gmapping}-style A*-plus-local-planner stack --- follows
the publicly available reference implementation
\textbf{Frontier\_Exploration} by Hanwen Cao,\footnote{\url{https://github.com/HanwenCao/Frontier_Exploration}}
which the project adapts and re-implements in Python for the Webots ROSbot. The
reference is cited explicitly to make the provenance of the algorithmic pattern
transparent and to avoid any appearance of plagiarism. The mapping, perception,
depth-auxiliary floating-wall layer, costmap centring, mission state machine,
observability, and per-maze configuration described above are original work
built for this Modularbeit; the public repository is credited as the source of
the underlying exploration-and-planning pattern rather than of the specific
code.

% ---------------------------------------------------------------------------
\section{Conclusion}

The approach delivers the mission --- reach blue, then yellow, in minimum time
without wall contact or poison traversal --- with a single, regression-tested
navigation algorithm shared across all five maze worlds and specialised to each
only through declarative configuration. Robustness to the mazes' lidar-blind
hazards is provided by a persistent depth-obstacle memory; tight-gap safety is
provided by footprint-accurate inflation, a centre-seeking costmap, and a
drift-immune DWA local planner; and correctness is safeguarded by a structured
event log and Webots-free regression gates. The common-controller design keeps
the algorithm auditable once while remaining adaptable to any new maze by
editing configuration alone.

\end{document}
